\documentclass{article}

% packages
\usepackage{amsmath, amsthm, thmtools, amsfonts, amssymb, luacode, catchfile, tikzducks, hyperref, ifthen}
\ifcsname c@kobocompile\endcsname
	\usepackage[a5paper, total={1072pt, 1448pt}, margin=10pt, includeheadfoot]{geometry} % set page margins
\else
	\usepackage[a4paper, margin=50pt, includeheadfoot]{geometry}
\fi
\usepackage[shortlabels]{enumitem}
\usepackage[skip=3pt, indent=0pt]{parskip}

% language
\usepackage[bidi=basic, layout=tabular, provide=*]{babel}
\ifcsname c@english\endcsname
	\babelprovide[main, import]{english}
\else
	\babelprovide[main, import]{hebrew}
	\babelprovide{rl}
\fi
%\babelfont{rm}{Libertinus Serif}
\babelfont{rm}[Renderer=Harfbuzz]{Libertinus Serif}
\babelfont{sf}{Libertinus Sans}
\babelfont{tt}{Libertinus Mono}

% style
\AddToHook{cmd/section/before}{\clearpage}	% Add line break before section
\linespread{1.3}
\setcounter{secnumdepth}{0}		% Remove default number tags from sections, this won't do well with theorems
\AtBeginDocument{\setlength{\belowdisplayskip}{3pt}}
\AtBeginDocument{\setlength{\abovedisplayskip}{3pt}}
\graphicspath{ {../images/} }

% operators
\DeclareMathOperator\cis{cis}
\DeclareMathOperator\Sp{Sp}
\DeclareMathOperator\tr{tr}
\DeclareMathOperator\im{Im}
\DeclareMathOperator\re{Re}
\DeclareMathOperator\diag{diag}
\DeclareMathOperator*\lowlim{\underline{lim}}
\DeclareMathOperator*\uplim{\overline{lim}}
\DeclareMathOperator\rng{rng}
\DeclareMathOperator\Sym{Sym}
\DeclareMathOperator\Arg{Arg}
\DeclareMathOperator\Log{Log}
\DeclareMathOperator\dom{dom}
\DeclareMathOperator\supp{Supp}
\DeclareMathOperator\var{Var}
\DeclareMathOperator\cov{Cov}

% commands
%\renewcommand\qedsymbol{\textbf{מש''ל}}
%\renewcommand\qedsymbol{\fbox{\emoji{lizard}}}
\newcommand{\Aa}[0]{\mathcal{A}}
\newcommand{\Bb}[0]{\mathcal{B}}
\newcommand{\CC}[0]{\mathbb{C}}
\newcommand{\Cc}[0]{\mathcal{C}}
\newcommand{\EE}[0]{\mathbb{E}}
\newcommand{\FF}[0]{\mathbb{F}}
\newcommand{\Ff}[0]{\mathcal{F}}
\newcommand{\Ii}[0]{\mathcal{I}}
\newcommand{\Gg}[0]{\mathcal{G}}
\newcommand{\Ll}[0]{\mathcal{L}}
\newcommand{\Mm}[0]{\mathcal{M}}
\newcommand{\NN}[0]{\mathbb{N}}
\newcommand{\Nn}[0]{\mathcal{N}}
\newcommand{\PP}[0]{\mathbb{P}}
\newcommand{\Pp}[0]{\mathcal{P}}
\newcommand{\QQ}[0]{\mathbb{Q}}
\newcommand{\RR}[0]{\mathbb{R}}
\newcommand{\Rr}[0]{\mathcal{R}}
\newcommand{\Ss}[0]{\mathcal{S}}
\newcommand{\TT}[0]{\mathbb{T}}
\newcommand{\Uu}[0]{\mathcal{U}}
\newcommand{\Vv}[0]{\mathcal{V}}
\newcommand{\Ww}[0]{\mathcal{W}}
\newcommand{\ZZ}[0]{\mathbb{Z}}
\newcommand{\acts}[0]{\circlearrowright}
\newcommand{\explain}[2] {
	\begin{flalign*}
		 && \text{#2} && \text{#1}
	\end{flalign*}
}
\newcommand{\maketitleprint}[0]{ \begin{center}
	%\begin{tikzpicture}[scale=3]
	%	\duck[graduate=gray!20!black, tassel=red!70!black]
	%\end{tikzpicture}	
	\includegraphics[width=6cm]{cover}
\end{center}
}

% theorem commands
\newtheoremstyle{c_remark}
	{}	% Space above
	{}	% Space below
	{}% Body font
	{}	% Indent amount
	{\bfseries}	% Theorem head font
	{}	% Punctuation after theorem head
	{.5em}	% Space after theorem head
	{\thmname{#1}\thmnumber{ #2}\thmnote{ \normalfont{\text{(#3)}}}}	% head content
\newtheoremstyle{c_definition}
	{3pt}	% Space above
	{3pt}	% Space below
	{}% Body font
	{}	% Indent amount
	{\bfseries}	% Theorem head font
	{}	% Punctuation after theorem head
	{.5em}	% Space after theorem head
	{\thmname{#1}\thmnumber{ #2}\thmnote{ \normalfont{\text{(#3)}}}}	% head content
\newtheoremstyle{c_plain}
	{3pt}	% Space above
	{3pt}	% Space below
	{\itshape}% Body font
	{}	% Indent amount
	{\bfseries}	% Theorem head font
	{}	% Punctuation after theorem head
	{.5em}	% Space after theorem head
	{\thmname{#1}\thmnumber{ #2}\thmnote{ \text{(#3)}}}	% head content

\ifcsname c@english\endcsname
	\theoremstyle{plain}
	\newtheorem{theorem}{Theorem}[section]
	\newtheorem{lemma}[theorem]{Lemma}
	\newtheorem{proposition}[theorem]{Proposition}
	\newtheorem*{proposition*}{Proposition}
	%\newtheorem{corollary}[theorem]{אין חלופה עברית}

	\theoremstyle{definition}
	\newtheorem{definition}[theorem]{Definition}
	\newtheorem*{definition*}{Definition}
	\newtheorem{example}{Example}[section]
	\newtheorem{exercise}{Exercise}[section]

	\theoremstyle{remark}
	\newtheorem*{remark}{Remark}
	\newtheorem*{solution}{Solution}
	\newtheorem{conclusion}[theorem]{Conclusion}
	\newtheorem{notation}[theorem]{Notation}
\else
	\theoremstyle{c_plain}
	\newtheorem{theorem}{משפט}[section]
	\newtheorem{lemma}[theorem]{למה}
	\newtheorem{proposition}[theorem]{טענה}
	\newtheorem*{proposition*}{טענה}
	%\newtheorem{corollary}[theorem]{אין חלופה עברית}

	\theoremstyle{c_definition}
	\newtheorem{definition}[theorem]{הגדרה}
	\newtheorem*{definition*}{הגדרה}
	\newtheorem{example}{דוגמה}[section]
	\newtheorem{exercise}{תרגיל}[section]

	\theoremstyle{c_remark}
	\newtheorem*{remark}{הערה}
	\newtheorem*{solution}{פתרון}
	\newtheorem{conclusion}[theorem]{מסקנה}
	\newtheorem{notation}[theorem]{סימון}
\fi

% Questions related commands
\newcounter{question}
\setcounter{question}{1}
\newcounter{sub_question}
\setcounter{sub_question}{1}

\ifcsname c@english\endcsname
	\newcommand{\question}[1][0]{
		\ifthenelse{#1 = 0}{}{\setcounter{question}{#1}}
		\section{Question \arabic{question}}
		\addtocounter{question}{1}
		\setcounter{sub_question}{1}
	}

	\newcommand{\subquestion}[1][0]{
		\ifthenelse{#1 = 0}{}{\setcounter{sub_question}{#1}}
		\subsection{Part \alph{sub_question}}
		\addtocounter{sub_question}{1}
	}
\else
	\newcommand{\question}[1][0]{
		\ifthenelse{#1 = 0}{}{\setcounter{question}{#1}}
		\section{שאלה \arabic{question}}
		\addtocounter{question}{1}
		\setcounter{sub_question}{1}
	}

	\newcommand{\subquestion}[1][0]{
		\ifthenelse{#1 = 0}{}{\setcounter{sub_question}{#1}}
		\subsection{סעיף \localecounter{letters.gershayim}{sub_question}}
		\addtocounter{sub_question}{1}
	}
\fi

% import lua and start of document
\directlua{common = require ('../common')}

\GetEnv{AUTHOR}

% headers
\author{\AUTHOR}
\date\today

\title{פתרון מטלה 2 --- גאומטריה דיפרנציאלית אלמנטרית, 80560}

\begin{document}
\maketitle
\maketitleprint[purple]

\question{}
יהי $(E, V, t)$ מרחב אפיני ממימד $n$ מעל $\FF$.

\subquestion{}
נראה שהאוסף ${\{ P_i \}}_{i \in I} \subseteq E$ בלתי־תלוי אפינית אם ורק אם לכל $i \in I$ האוסף ${\{ P_j - P_i \}}_{i \ne j \in I} \subseteq V$ בלתי־תלויה לינארית מעל $V$.
\begin{proof}
	נניח ש־$\{ P_i \}$ בלתי־תלוי אפינית ויהי $i \in I$.
	נניח בשלילה גם ש־$B = {\{ P_j - P_i \}}_{i \ne j \in I}$ היא תלויה לינארית, ובפרט נניח שמתקיים,
	\[
		P_0 - P_i
		\in \Sp\{ P_j - P_i \mid j \in I \setminus \{0, i\} \}
	\]
	כאשר $i \ne 0 \in I$ בלי הגבלת הכלליות.
	נזכור שמתקיים $P_0 \notin \langle P_j \mid j \ne i \rangle = P_i + W$ מאי־התלות האפינית, כאשר $W \le V$.
	אבל $P_0 - P_i \in W$ מהטענה שראינו זה עתה, ולכן בפרט $P_0 \notin P_i + W$ בסתירה, ולכן הקבוצה לא תלויה לינארית.

	נניח עתה את הכיוון ההפוך.
	נראה שמתקיים $P_0 \notin \langle P_j \mid j \ne 0 \rangle = P_i + W$ עבור $W \le V$.
	נניח בשלילה שאכן $P_0 \in P_i + W$, כלומר קיים $u \in W$ כך שמתקיים $P_0 = P_i + u$, כלומר $P_0 - P_i = u$.
	אבל הנחנו ש־$u \in \{ P_j - P_i \mid j \ne i \}$ היא בלתי־תלויה לינארית, וזו סתירה.
\end{proof}

\subquestion{}
נראה ש־${\{ P_i \}}_{0 \le i \le k}$ בלתי־תלויה אפינית אם ורק אם $\dim{\langle P_i \rangle}_{i \le k} = k$.
\begin{proof}
	נניח שהקבוצה בלתי־תלויה אפינית, ולכן $\langle P_i \mid i \le k \rangle = P_0 + W$ עבור $W \le V$ ומהסעיף הקודם $\dim W = k$.

	בכיוון ההפוך נניח ש־$\dim W = k$.
	אבל בקבוצה $\{ P_i - P_0 \mid 0 < i \le k \} = B$ יש בדיוק $k$ איברים ולכן $W = \Sp B$ מקיים את תנאי סעיף א' ונקבל ש־$\{ P_i \mid i \le k \}$ בלתי־תלויה אפינית.
\end{proof}

\subquestion{}
נראה ש־${\{ P_i \}}_{0 \le i \le n}$ בסיס אפיני של $E$ אם ורק אם לכל $i \in I$ מתקיים,
\[
	(P_i, (P_0 - P_i, \ldots, P_n - P_i))
\]
היא מערכת יחוס של $E$.
\begin{proof}
	נניח ש־$(P_0, \ldots, P_n)$ בסיס אפיני ולכן $\langle P_0, \ldots, P_n \rangle = P + V$ עבור $P_i \in E$ כלשהי (נובע מסעיף ב'), עבור $i \in I$ נתון.
	נסיק מסעיף א' ש־$V = \Sp\{ P_0 - P_i, \ldots, P_n - P_i \}$\ ולכן נקבל מערכת יחוס.

	בכיוון ההפוך מהגדרת מערכות יחוס נקבל ש־$\Sp\{ P_0 - P_i, \ldots, P_n - P_i \} = V$ ולכן $P_i + V = E$, ובפרט מהנתון אודות כל $i$ נקבל מסעיף א' ש־$\{ P_i \}$ בסיס אפיני של $E$.
\end{proof}

\question{}
יהיו $(E, V), (F, U), (G, W)$ שלושה מרחבים אפיניים מעל השדה $\FF$.
תהיינה $f : E \to F, g : F \to G$ העתקות אפיניות. \\
נראה ש־$g \circ f : E \to G$ העתקה אפינית המקיימת $d(g \circ f) = (dg) \circ (df)$.
\begin{proof}
	נסמן $g \circ f$, נתחיל להראות ש־$dh$ מוגדרת היטב, כלומר שלכל $u \in V$ מתקיים $h(P + u) - h(P) = h(Q + u) - h(Q)$.
	\[
		h(P + u) - h(P)
		= g(f(P + u)) - h(P)
		= g(f(P) + df(u)) - h(P)
		= g(f(P)) + dg(df(u)) - h(P)
		= dg(df(u)) 
	\]
	הביטוי לא תלוי ב־$P$ ולכן הפונקציה מוגדרת היטב ונבחין כי היא גם לינארית כהרכבת העתקות לינאריות.
	נסיק בהתאם ש־$h$ היא העתקה אפינית, וישירות מהגדרת $dh$ נקבל מהחישוב שעשינו ש־$dh = (dg) \circ (df)$.
\end{proof}

\question{}
תהיינה $f^1, \ldots, f^m : \FF^n \to \FF$ ותהי $f : \FF^n \to \FF^m$ המוגדרת על־ידי $f(x) = {(f^1(x), \ldots, f^m(x))}^t$.

\subquestion{}
נוכיח ש־$f$ לינארית אם ורק אם $f^1, \ldots, f^m$ העתקות לינאריות.
\begin{proof}
	נגדיר את ההעתקה הלינארית $H_k : \FF^m \to \FF$ המוגדרת על־ידי $H_k(x_1, \ldots, x_m) = x_k$.
	זוהי העתקת הטלה וידוע שהיא לינארית.
	נבחין כי מתקיים $H_k \circ f \equiv f^k$ לכל $k$, אבל $f$ לינארית מהנחה ו־$H_k$ לינארית ולכן ההרכבה שלהן היא לינארית, כלומר $f^k$ לינארית לכל $1 \le k \le n$.

	נניח ש־$f^1, \ldots, f^n$ לינאריות.
	נראה ש־$f$ לינארית ישירות מהגדרה.
	נניח ש־$x, y \in \FF^n$ וכן ש־$\alpha, \beta \in \FF$, אז מתקיים,
	\[
		f(\alpha x + \beta y)
		= \begin{pmatrix} f^1(\alpha x + \beta y) \\ \vdots \\ f^m(\alpha x + \beta y) \end{pmatrix} 
		= \begin{pmatrix} \alpha f^1(x) + \beta f^1(y) \\ \vdots \\ \alpha f^m(x) + \beta f^m(y) \end{pmatrix} 
		= \alpha \begin{pmatrix} f^1(x) \\ \vdots \\ f^m(x) \end{pmatrix} + \beta \begin{pmatrix} f^1(y) \\ \vdots \\ f^m(y) \end{pmatrix} 
		= \alpha f(x) + \beta f(y)
	\]
	ולכן $f$ לינארית.
\end{proof}

\subquestion{}
נוכיח ש־$f$ העתקה אפינית אם ורק אם $f^1, \ldots, f^m$ העתקות אפיניות.
\begin{proof}
	נניח ש־$f$ אפינית.
	
	TODO
\end{proof}

\end{document}
